\section*{Problématique générale}


Ce manuscrit aborde le problème de la structuration automatique de données. La question qui le traverse est la suivante : comment extraire, à partir d'observations locales, imparfaites et parfois bruitées, une structure globale qui soit robuste, interprétable et mathématiquement contrôlée ? Cette question apparaît dans plusieurs contextes : classification non supervisée de nuages de points, nettoyage de données LiDAR, segmentation d'objets dans des scènes 3D ou 4D, détection de communautés dans des graphes, ou extraction de réseaux de fissures dans des images. Ces situations semblent différentes. Elles posent pourtant un problème commun : il faut passer d'informations locales, souvent ambiguës, à une organisation globale des données.

Le point de départ du manuscrit est le clustering, ou classification non supervisée, dans un nuage fini de points
\[
X = \{x_1,\ldots,x_n\}\subset \mathbb{R}^p.
\]
Ce cadre euclidien est volontairement restreint. Il permet cependant de faire apparaître les phénomènes étudiés dans toute la suite : la hiérarchie, la densité, la connexité, le bruit et la percolation. Le modèle statistique de référence est celui de Hartigan. Dans ce modèle, les clusters ne sont pas définis par des centres, mais comme les composantes connexes des ensembles de niveau d'une densité $f$. Pour un niveau $\lambda>0$, l'ensemble de forte densité est
\[
L_f(\lambda)=\{x\in \mathbb{R}^p\mid f(x)\geq \lambda\},
\]
et les amas de forte densité sont les composantes connexes de cet ensemble.

En pratique, la densité $f$ est inconnue. Il faut donc l'estimer. L'estimateur qui sert de fil conducteur au manuscrit est l'estimateur des $K$ plus proches voisins, ou $K$-NN. En paramétrant les niveaux de densité par un rayon $r>0$, les ensembles de niveau s'écrivent
\[
L_K(r)=\{y\in \mathbb{R}^p\mid |B(y,r)\cap X|\geq K\}.
\]
Un point $y$ appartient donc à $L_K(r)$ lorsque la boule centrée en $y$ et de rayon $r$ contient au moins $K$ points du nuage. Pour $K=1$, ces ensembles de niveau sont simplement des unions de boules, et leurs composantes connexes sont décrites par un graphe géométrique. Ce cas correspond exactement à l'algorithme du Single-Linkage. La question centrale devient alors : que devient cette connexité lorsque $K\geq 2$ ?

La réponse proposée dans ce manuscrit est que les graphes ordinaires ne suffisent plus. Lorsque $K\geq 2$, la définition même de $L_K(r)$ fait intervenir des intersections simultanées de $K$ boules. Les objets élémentaires ne sont donc plus seulement des arêtes entre deux points, mais des interactions entre plusieurs points. Il faut alors remplacer les graphes par des complexes simpliciaux, en particulier le complexe de \v{C}ech, et remplacer les composantes connexes usuelles par une notion de connexité d'ordre supérieur. Cette idée conduit à l'algorithme HGP-Clusterer, pour Hypergraphe-Percol'.

Le second fil conducteur est la percolation. Dans un graphe géométrique, lorsque le rayon de connexion augmente, de petites composantes apparaissent, croissent, puis fusionnent. En présence de bruit, une composante géante peut apparaître dans le fond et relier des régions qui devraient rester séparées. Ce phénomène explique l'effet de chaînage du Single-Linkage, mais aussi certaines limites de ses robustifications modernes, comme le Robust Single-Linkage, DBSCAN et HDBSCAN. La percolation fournit ainsi un langage pour comprendre les échecs des méthodes par graphes et pour comparer plus finement les méthodes fondées sur des interactions d'ordre supérieur.

Le manuscrit est organisé en trois parties. La première revient sur le Single-Linkage, ses fondements et ses limites. La deuxième construit la généralisation par interactions d'ordre supérieur et l'analyse par percolation. La troisième montre que ces idées dépassent le cadre des nuages de points euclidiens, en les appliquant à des graphes pondérés signés et à des images.

\section*{Première partie : du Single-Linkage à ses fondements}

La première partie a pour objet l'algorithme du Single-Linkage. Elle ne se limite pas à une présentation algorithmique : elle en propose une lecture unifiée selon trois points de vue. Le premier est heuristique. Le Single-Linkage construit une hiérarchie en fusionnant, à chaque étape, les deux clusters les plus proches pour la distance du saut minimum. Dans un espace métrique fini, cette procédure est équivalente à la construction d'un arbre minimum couvrant. Toute l'information hiérarchique est contenue dans cet arbre, que l'on élague ensuite à différents niveaux de résolution. Dans les espaces euclidiens de basse dimension, cette construction peut être accélérée grâce à des graphes géométriques plus parcimonieux, notamment le graphe de Gabriel et la triangulation de Delaunay.

Le deuxième point de vue est géométrique. On peut définir, pour chaque rayon $r$, un graphe géométrique qui relie deux points lorsque leur distance est inférieure ou égale à $2r$. Lorsque $r$ varie, les composantes connexes de ce graphe apparaissent et fusionnent. Le Single-Linkage est exactement la hiérarchie obtenue en suivant cette évolution. Cette lecture est importante, car elle rapproche le clustering hiérarchique de l'analyse persistante : au lieu de produire une seule partition, l'algorithme décrit l'évolution de la connexité à toutes les échelles.

Le troisième point de vue est statistique. Dans le modèle de Hartigan, les clusters sont les composantes connexes des ensembles de niveau d'une densité. Lorsque cette densité est estimée par l'estimateur $1$-NN, les ensembles de niveau sont des unions de boules centrées sur les points du nuage. Les composantes de ces unions de boules coïncident avec les composantes du graphe géométrique, et donc avec l'arbre produit par le Single-Linkage. On obtient ainsi une identité entre trois objets : l'arbre minimum couvrant, la persistance des graphes géométriques et les amas discrets de forte densité associés à l'estimateur $1$-NN.

Cette triple lecture explique la place particulière du Single-Linkage. L'algorithme n'est pas seulement une heuristique historique. Il est aussi la solution exacte d'un modèle statistique précis, dans le cas particulier $K=1$. Le manuscrit rappelle également le point de vue axiomatique. Le théorème d'impossibilité de Kleinberg montre qu'aucune méthode de clustering produisant une partition unique ne peut satisfaire simultanément plusieurs propriétés naturelles. En demandant une hiérarchie plutôt qu'une partition unique, cette obstruction peut être contournée. Les travaux de Carlsson et Mémoli caractérisent alors le Single-Linkage comme une méthode canonique de projection des espaces métriques finis vers les espaces ultramétriques. Ce détour axiomatique n'est pas le centre du travail, mais il montre que le Single-Linkage réapparaît dans plusieurs cadres mathématiques distincts.

Les limites du Single-Linkage sont ensuite étudiées. Son défaut principal est l'effet de chaînage. Une suite de points isolés peut suffire à relier deux amas que l'on voudrait séparer. Dans le modèle de Hartigan, ce problème se traduit par une inconsistance en dimension $p\geq 2$. Le phénomène n'est pas seulement algorithmique : il est lié à la percolation. Lorsque le rayon augmente, le bruit peut former des ponts de connexité avant que les amas de forte densité ne soient correctement récupérés. En dimension un, cette difficulté ne se manifeste pas de la même façon ; en dimension supérieure, elle devient structurelle.

Les méthodes modernes comme Robust Single-Linkage, DBSCAN et HDBSCAN cherchent à corriger ce problème en introduisant une densité locale fondée sur les $K$ plus proches voisins. Le Robust Single-Linkage impose que les points connectés soient eux-mêmes localement denses. DBSCAN ajoute la notion de points accessibles, ce qui permet de mieux gérer les frontières des clusters. HDBSCAN construit une hiérarchie, élaguée par un seuil de taille minimale, puis sélectionne les clusters au moyen du critère d'excès de masse. Ces méthodes constituent une amélioration pratique importante. Toutefois, le manuscrit formule une critique précise : elles filtrent les points avec une condition d'ordre $K$, mais elles propagent encore la connexité par des arêtes ordinaires. La densité est donc d'ordre $K$, tandis que la connexité reste binaire.

Cette asymétrie explique le paradoxe du choix de $K$ dans HDBSCAN. En théorie, augmenter $K$ devrait lisser l'estimateur de densité et améliorer sa consistance statistique. En pratique, des valeurs trop grandes dégradent la géométrie des clusters. Le problème ne vient pas seulement d'un lissage excessif : il vient du fait que le graphe construit par ces méthodes ne représente plus correctement les ensembles de niveau $K$-NN. Pour corriger cela, il ne suffit donc pas de modifier une distance ou un critère d'élagage. Il faut modifier la notion même de connexité.

La première partie se termine par une utilisation plus pratique de la hiérarchie. Un arbre ultramétrique peut servir d'espace de résolution. Des problèmes difficiles dans l'espace euclidien deviennent plus simples lorsqu'ils sont posés sur une hiérarchie. Le manuscrit montre comment explorer un arbre hiérarchique en remplaçant le critère statistique d'HDBSCAN par des a priori géométriques propres au problème. Deux applications sont présentées. La première concerne la détection d'anomalies dans un nuage LiDAR 3D, avec un modèle théorique de chantier naval utilisé comme information préalable. La seconde concerne la segmentation panoptique 4D sur SemanticKITTI : les objets sont segmentés dans chaque trame, puis suivis dans le temps à l'aide du transport optimal déséquilibré régularisé par l'entropie. Ces applications montrent que la hiérarchie n'est pas seulement une sortie à visualiser. Elle peut aussi devenir un espace dans lequel on injecte des contraintes du problème.

\section*{Deuxième partie : généralisation du Single-Linkage par interactions d'ordre supérieur}

La deuxième partie constitue le coeur théorique du manuscrit. Son objectif est de généraliser le Single-Linkage aux estimateurs $K$-NN, tout en conservant l'exactitude obtenue dans le cas $K=1$. Pour cela, le manuscrit repart de la définition des ensembles de niveau
\[
L_K(r)=\{y\in \mathbb{R}^p\mid |B(y,r)\cap X|\geq K\}.
\]
Dire que $y$ appartient à $L_K(r)$ signifie que $y$ appartient simultanément à au moins $K$ boules centrées sur des points de $X$. La bonne structure n'est donc plus le graphe géométrique, qui encode des relations deux à deux, mais le complexe de \v{C}ech
\[
\check C(X,r)=\left\{\sigma\subset X\;\middle|\;\bigcap_{x\in\sigma}B(x,r)\neq\varnothing\right\}.
\]
Un simplexe du complexe encode l'existence d'une intersection commune de boules. Les $(K-1)$-simplexes représentent donc les régions où $K$ points du nuage peuvent être vus simultanément depuis une même boule de rayon $r$.

Le manuscrit définit alors les $K$-polyèdres. On construit un graphe auxiliaire $\Gamma_K(X,r)$ dont les sommets sont les $(K-1)$-simplexes du complexe de \v{C}ech. Deux tels simplexes $\sigma$ et $\tau$ sont adjacents lorsque leur union $\sigma\cup\tau$ est encore un simplexe du complexe. Les composantes connexes de ce graphe auxiliaire définissent les $K$-polyèdres. Pour $K=1$, on retrouve les composantes connexes ordinaires du graphe géométrique, donc le Single-Linkage. Pour $K=2$, ce ne sont plus les points qui percolent directement, mais les arêtes, recollées entre elles par des triangles. Plus généralement, la connexité se propage par interactions d'ordre supérieur.

L'algorithme obtenu est appelé HGP-Clusterer. Il observe l'évolution des $K$-polyèdres lorsque $r$ varie. Le résultat principal est l'identité
\[
\theta^{\mathrm{HGP}}_K(r)=H^{\mathrm{discrets}}_{\widehat f_K}(r),
\]
c'est-à-dire que les $K$-polyèdres du complexe de \v{C}ech coïncident, niveau par niveau, avec les amas discrets de forte densité associés à l'estimateur $K$-NN. Cette identité généralise exactement celle obtenue pour le Single-Linkage lorsque $K=1$.

La preuve repose sur les régions témoins. À tout $(K-1)$-simplexe $\sigma$, on associe
\[
T_r(\sigma)=\bigcap_{x\in\sigma}B(x,r).
\]
Le niveau $L_K(r)$ est l'union de ces régions témoins. Le graphe $\Gamma_K(X,r)$ est le graphe d'intersection de cette famille de régions. Comme ces régions sont convexes et connexes, les composantes connexes de leur union correspondent aux composantes connexes du graphe d'intersection. On obtient ainsi la correspondance entre composantes de $L_K(r)$ et $K$-polyèdres. Le passage au discret se fait ensuite en prenant les points de $X$ couverts par ces composantes.

Un point important est que, pour $K\geq 2$, les amas discrets de forte densité ne forment pas nécessairement une partition des points. Un même point peut participer à plusieurs interactions d'ordre supérieur et appartenir à plusieurs $K$-polyèdres. Ce recouvrement n'est pas une anomalie de la méthode. Il correspond à l'information géométrique présente dans les données. Forcer immédiatement une partition ferait perdre cette information. Lorsque l'application impose une partition stricte, le manuscrit propose un post-traitement par vote pondéré.

L'exactitude topologique ne suffit toutefois pas à caractériser les performances statistiques. La question suivante est donc étudiée : quelle fraction d'un amas de forte densité peut-on récupérer avant que le bruit ne provoque une fusion parasite ? Cette question conduit à la théorie de la percolation continue. Le manuscrit compare trois familles de composantes : les coeurs du Robust Single-Linkage, les composantes de DBSCAN/HDBSCAN, et les $K$-polyèdres de HGP-Clusterer. À chacune est associée une probabilité de percolation
\[
\lambda\longmapsto \Theta^{cc}_{K,p}(\lambda),
\qquad cc\in\{\mathrm{core},\mathrm{dbscan},\mathrm{poly}\}.
\]
Cette fonction mesure la probabilité qu'un point appartienne à une composante macroscopique. Elle permet de comparer les méthodes non pas seulement par un score empirique, mais par la manière dont leurs composantes apparaissent et fusionnent.

En dimension $p\geq 2$, à $K$ fixé, il n'est pas raisonnable d'espérer récupérer intégralement deux amas denses avant que le fond ne percole. L'objectif devient de mesurer la fraction asymptotiquement récupérable. Le manuscrit introduit pour cela une vitesse de percolation. Plus cette vitesse est favorable, plus la méthode récupère rapidement une grande partie de l'amas dense avant que le bruit ne forme une composante géante. Les simulations en petites dimensions montrent que les $K$-polyèdres ont une meilleure vitesse de percolation que le $K$-Robust Single-Linkage et que $K$-DBSCAN. L'interprétation est directe : il ne suffit pas de demander aux points d'être denses ; il faut connecter correctement les régions où la densité $K$-NN est forte.

Le manuscrit étudie aussi le régime $K\to +\infty$. Dans la bonne fenêtre de normalisation, les niveaux de forte densité $K$-NN sont remplacés par les excursions d'un champ gaussien stationnaire. Cette limite permet de distinguer les obstructions propres à chaque méthode. Pour les coeurs, l'échec peut venir d'une fluctuation ponctuelle. Pour DBSCAN/HDBSCAN, il suffit d'une obstruction de rayon fini. Pour les $K$-polyèdres, l'échec est lié à une obstruction géométrique de connexion. Cette lecture explique pourquoi les interactions d'ordre supérieur améliorent la consistance fractionnelle.

Reste la question algorithmique. Construire tout le complexe de \v{C}ech est prohibitif, car le nombre de simplexes peut croître très vite. Le manuscrit cherche donc l'analogue, à l'ordre $K$, de l'arbre minimum couvrant du Single-Linkage. L'idée est de conserver seulement les simplexes qui provoquent effectivement des fusions dans la hiérarchie. Ces simplexes sont appelés $K$-séparants. Le manuscrit montre qu'un simplexe $K$-séparant satisfait nécessairement une condition de Gabriel, c'est-à-dire une condition de vacuité géométrique. Les fusions utiles sont donc contenues dans une structure plus mince que le complexe complet.

Cette structure est reliée à la mosaïque de Delaunay d'ordre $K$. De même que l'arbre minimum couvrant du Single-Linkage est contenu dans le graphe de Gabriel, lui-même lié à la triangulation de Delaunay, le $K$-arbre minimum couvrant de HGP-Clusterer s'appuie sur les simplexes de Gabriel et sur la mosaïque de Delaunay d'ordre supérieur. Le triptyque du Single-Linkage est ainsi reconstruit : un modèle statistique, une analyse persistante géométrique, et un objet combinatoire plus mince pour le calcul.

Le chapitre pratique précise ensuite comment utiliser HGP-Clusterer. Lorsque l'on veut obtenir une partition stricte, un vote pondéré attribue chaque point à un cluster unique à partir des faces incidentes et de leurs scores de densité. Pour $K=2$, une mise en oeuvre optimisée s'appuie sur la triangulation de Delaunay standard. Des expériences sont présentées sur des données synthétiques et sur le jeu de données des huiles d'olive italiennes, décrit par huit concentrations en acides gras. Ces résultats illustrent l'intérêt de l'ordre $2$ pour séparer des structures que les méthodes fondées uniquement sur des graphes représentent moins finement. Lorsque la dimension est trop grande, ou lorsque les données ne sont pas euclidiennes, le complexe de Vietoris--Rips fournit une approximation naturelle : les $K$-polyèdres deviennent alors des composantes de percolation de cliques dans un graphe seuillé.

\section*{Troisième partie : vers des données davantage structurées}

La troisième partie montre que les deux idées principales du manuscrit ne se limitent pas aux nuages de points euclidiens. Les interactions d'ordre supérieur et la percolation réapparaissent dans deux cadres différents : les graphes pondérés signés et les images.

Le premier cadre est celui de la détection de communautés. Les données sont des graphes non orientés, pondérés, signés, éventuellement munis d'attributs sur les sommets. Les sommets appartiennent à des communautés cachées, et les arêtes observées donnent des informations bruitées sur ces appartenances. Une arête positive favorise l'égalité des étiquettes ; une arête négative favorise leur différence. Cette situation peut produire de la frustration : certaines contraintes locales ne peuvent pas être satisfaites simultanément. Le triangle pleinement répulsif est le plus petit exemple de cette obstruction.

Le manuscrit construit un cadre bayésien pour ce problème. Il définit une loi a priori sur les communautés, un canal d'observation pour les données, puis une loi postérieure sur les configurations possibles. La performance d'un algorithme est formulée en termes de recouvrement : battre un tirage au hasard, obtenir un recouvrement partiel, presque parfait ou parfait. Cette formalisation permet de traiter la détection de communautés comme un problème d'inférence probabiliste, et non seulement comme une optimisation de score.

Pour explorer la loi postérieure, le manuscrit introduit des dynamiques de clusters inspirées de Swendsen--Wang. Ces dynamiques modifient des groupes de variables plutôt que des étiquettes une par une. Elles sont adaptées aux mesures de Gibbs et peuvent être généralisées aux interactions d'ordre supérieur. Le cas triangulaire est particulièrement utile pour traiter la frustration. Sur un triangle attractif, il est possible de geler simultanément plusieurs contraintes. Sur un triangle répulsif, un gel complet introduirait une contradiction ; les règles locales doivent donc préserver la réversibilité tout en évitant de figer une structure impossible à satisfaire.

La percolation réapparaît alors comme condition de détectabilité. Dans le modèle de Sankararaman--Baccelli, le manuscrit retrouve les bornes connues en appliquant une dynamique de Swendsen--Wang arête par arête. Il montre ensuite qu'une dynamique triangulaire permet d'obtenir des bornes plus fines dans le cas de la grille triangulaire. L'idée générale est analogue à celle développée pour HGP-Clusterer : une interaction d'ordre supérieur peut réduire certaines ambiguïtés produites par les relations binaires et fournir une meilleure lecture des limites de récupération. Le chapitre se termine par une ouverture vers l'assemblage d'haplotypes, que l'on peut formuler comme un problème de détection de communautés sur graphe signé.

Le second cadre est celui du traitement d'image, avec l'extraction de réseaux de fissures. Les méthodes d'apprentissage profond dominent aujourd'hui ce domaine, mais elles demandent des annotations et peuvent se dégrader lorsque les capteurs, les matériaux ou les résolutions changent. Le manuscrit propose une approche sans apprentissage, fondée sur des modèles géométriques explicites. Le point de départ est le filtre de Frangi, qui détecte localement les structures tubulaires ou linéiques à partir des valeurs propres de la matrice hessienne de l'image. Ce filtre est rapide et interprétable, mais il reste local : il indique qu'un pixel ressemble à un morceau de fissure, sans dire s'il appartient au même réseau qu'un autre pixel voisin.

Le manuscrit introduit donc le graphe de Frangi. Les sommets sont des pixels candidats, et les arêtes mesurent la compatibilité entre deux morceaux possibles de fissure. Cette compatibilité tient compte de l'élongation, de la courbure transverse et de l'alignement des directions principales. La sortie visée n'est pas seulement un masque binaire, mais un squelette ou un graphe, donc un objet géométrique éditable. Ce choix est adapté à l'annotation experte : un graphe peut être corrigé, simplifié ou validé plus facilement qu'une carte de probabilités pixel par pixel.

L'approche s'étend naturellement à la multimodalité. Lorsque plusieurs images recalées sont disponibles, par exemple une image visible, une carte de profondeur ou une image thermique, les opérateurs différentiels peuvent être normalisés puis fusionnés. L'hypothèse est qu'une fissure physique produit une structure cohérente entre les modalités, tandis que certains artefacts sont plus spécifiques à une seule modalité. Les expériences sur FIND montrent l'intérêt de cette fusion : la profondeur et la profondeur filtrée apportent une information complémentaire à l'intensité visible.

Les interactions d'ordre supérieur sont utilisées dans les images bruitées ou fortement texturées. Sur VT-GraF, les cailloux et les textures granulaires produisent de nombreuses réponses locales allongées. Un graphe de Frangi ordinaire peut alors créer des branches parasites. En passant à $K=2$, les arêtes doivent appartenir à des triangles ; cette contrainte agit comme un filtre de densité locale et élimine une partie des réponses isolées. Les expériences sur Vaches Noires, FIND, VT-GraF et CrackForest montrent que la méthode peut être compétitive sans annotation, en particulier lorsque les données sortent de la distribution sur laquelle des modèles appris sont entraînés. Sur des données propres et proches de l'entraînement, les méthodes profondes restent souvent meilleures, ce qui est cohérent avec leur spécialisation.

Le chapitre souligne aussi les limites de cette approche. Les paramètres d'échelle restent importants, car ils doivent correspondre à la largeur attendue des fissures. Le rayon du graphe doit être choisi avec soin : trop petit, il fragmente le réseau ; trop grand, il relie artificiellement des structures voisines. Les seuils relatifs fonctionnent lorsque l'on sait que l'image contient une fissure, mais une utilisation industrielle demanderait aussi un seuil absolu pour rejeter les images sans fissure. Ces limites n'annulent pas l'intérêt de la méthode : elles précisent les conditions dans lesquelles le graphe de Frangi peut devenir un outil d'aide à l'annotation ou un module complémentaire à des méthodes apprises.

\section*{Conclusion et perspectives}

Le manuscrit développe une ligne directrice unique à partir du Single-Linkage. Il commence par montrer que cet algorithme possède trois lectures équivalentes dans le cadre euclidien : arbre minimum couvrant, persistance des graphes géométriques et modèle statistique de Hartigan pour l'estimateur $1$-NN. Cette équivalence explique l'intérêt du Single-Linkage, mais révèle aussi sa limite : la connexité binaire par arêtes est insuffisante dès que l'on veut utiliser correctement une densité $K$-NN avec $K\geq 2$.

La contribution centrale est HGP-Clusterer. En remplaçant le graphe géométrique par le complexe de \v{C}ech, et les composantes connexes ordinaires par les $K$-polyèdres, l'algorithme retrouve exactement les amas discrets de forte densité de l'estimateur $K$-NN. Il ne s'agit donc pas seulement d'une modification heuristique de HDBSCAN ou de DBSCAN. La construction restaure l'accord entre la définition statistique de la densité locale et la notion de connexité utilisée pour construire la hiérarchie.

La percolation fournit le cadre d'analyse des performances. Elle explique l'effet de chaînage, les fusions parasites dans le bruit, et la consistance seulement fractionnelle en dimension supérieure ou égale à deux. Elle permet aussi de comparer les méthodes par leurs courbes de percolation et par la fraction d'amas récupérable avant fusion. Dans cette lecture, les interactions d'ordre supérieur ne sont pas seulement plus expressives ; elles modifient la manière dont les structures apparaissent et se connectent.

Le manuscrit traite également la question du calcul. Les simplexes $K$-séparants, les conditions de Gabriel et la mosaïque de Delaunay d'ordre $K$ fournissent l'analogue d'ordre supérieur de l'arbre minimum couvrant. Pour $K=2$, des optimisations pratiques rendent l'approche utilisable, et des approximations par Vietoris--Rips permettent d'envisager des espaces de plus grande dimension ou non euclidiens.

Enfin, les applications aux graphes signés et aux images montrent que les idées développées dans le cadre euclidien ont une portée plus large. Dans les graphes, les interactions triangulaires aident à traiter la frustration et à affiner les bornes de recouvrement. Dans les images, le graphe puis l'hypergraphe de Frangi donnent une méthode géométrique sans apprentissage pour extraire des réseaux de fissures, avec une sortie éditable.

Plusieurs perspectives restent ouvertes. La première concerne l'axiomatique des recouvrements : si les $K$-polyèdres ne forment pas toujours des partitions, il faut mieux comprendre les bons espaces de sortie pour des hiérarchies de recouvrements. La deuxième concerne la consistance statistique en dimension $p\geq 2$, où la percolation impose des limites structurelles. La troisième concerne la réduction de dimension : la mosaïque de Delaunay d'ordre $K$ encode des relations locales plus riches qu'un graphe de voisinage, ce qui suggère un analogue d'ordre supérieur d'UMAP. Enfin, pour l'imagerie, les graphes de Frangi pourraient servir à produire des pré-annotations corrigées par des experts, ou être intégrés à des architectures apprises sous forme de contraintes géométriques explicites.

